\documentclass[11pt,a4paper]{ctexart}
\usepackage{geometry}
\geometry{left=1.8cm,right=1.8cm,top=1.5cm,bottom=2cm}
\usepackage{fancyhdr}
\usepackage{enumitem}
\usepackage{xcolor}
\usepackage{hyperref}
\hypersetup{colorlinks=true,linkcolor=blue,urlcolor=blue}

\fancyhf{}
\fancyhead[L]{generated\_eval\_set\_v4 抽样检查}
\fancyhead[R]{\thepage}
\pagestyle{fancy}

\title{\textbf{电力审查数据集 v4 — 抽样检查（30题）}}
\author{L1/L2/L3 各随机抽取10题 \quad 完整字段展示}
\date{\today}

\begin{document}
\maketitle
\thispagestyle{fancy}

\section*{说明}
从 generated\_eval\_set\_v4.json（300题）中随机抽取 L1/L2/L3 各10题，
展示全部字段，供人工质量检查。

v4 相对于 v3 新增字段：
\textbf{clause\_source}（条款溯源）、
\textbf{sub\_category}（细粒度技术子类别）、
\textbf{LoopJudge verdict/score/gap}（弱/强模型评分闭环）。

\vspace{6pt}

\section*{L1 — 直接型（参数检索，答案在单一规范条款中可直接检索）}

\noindent\textbf{L1-007} \hfill \textbf{参数检索} $|$ \textit{电压等级与范围}
\begin{description}[leftmargin=1.5em,style=nextline,font=\normalfont]
  \item[Q:] 根据GB 38755-2019，该标准适用于电压等级为多少及以上的电力系统？
  \item[A:] 220kV
  \item[KW:] GB 38755-2019、电压等级、220kV、电力系统、适用范围
  \item[标准:] GB 38755-2019 §3.2.5.2
  \item[条款溯源:] GB 38755-2019 第3.5.1条
  \item[主题:] 通用
  \item[评分方式:] auto\_keyword\_match
  \item[Rubric条款:] GB 38755-2019 第3.5.1条
  \item[Rubric判据:] 答案应为: 220kV

\end{description}
\vspace{4pt}
\hrule
\vspace{6pt}

\noindent\textbf{L1-035} \hfill \textbf{参数检索} $|$ \textit{数值阈值与限值}
\begin{description}[leftmargin=1.5em,style=nextline,font=\normalfont]
  \item[Q:] 电网任一点综合零序电抗不应大于综合正序电抗的多少倍？
  \item[A:] 电网任一点综合零序电抗不应大于综合正序电抗的3倍。
  \item[KW:] 综合零序电抗、综合正序电抗、不应大于、3倍
  \item[标准:] DL/T 5429-2009 §6.5.3
  \item[条款溯源:] DL/T 5429-2009 第6.5.3条
  \item[主题:] 通用
  \item[评分方式:] auto\_keyword\_match
  \item[Rubric条款:] DL/T 5429-2009 第6.5.3条
  \item[Rubric判据:] 答案应为: 电网任一点综合零序电抗不应大于综合正序电抗的3倍。

\end{description}
\vspace{4pt}
\hrule
\vspace{6pt}

\noindent\textbf{L1-012} \hfill \textbf{参数检索} $|$ \textit{数值阈值与限值}
\begin{description}[leftmargin=1.5em,style=nextline,font=\normalfont]
  \item[Q:] 根据GB 38755-2019，N-1原则下，其他元件的负荷应满足什么具体数值条件？
  \item[A:] 不过负荷（即不超过额定容量）
  \item[KW:] N-1原则、不过负荷、额定容量、GB 38755-2019、元件
  \item[标准:] GB 38755-2019 §2.3
  \item[条款溯源:] GB 38755-2019 第2.3条
  \item[主题:] N-1原则
  \item[评分方式:] auto\_keyword\_match
  \item[Rubric条款:] GB 38755-2019 第2.3条
  \item[Rubric判据:] 答案应为: 不过负荷（即不超过额定容量）

\end{description}
\vspace{4pt}
\hrule
\vspace{6pt}

\noindent\textbf{L1-099} \hfill \textbf{参数检索} $|$ \textit{数值阈值与限值}
\begin{description}[leftmargin=1.5em,style=nextline,font=\normalfont]
  \item[Q:] 根据DL/T 5429-2009，电力系统远景展望水平年一般取今后多少年？
  \item[A:] 10\textasciitilde{}15年
  \item[KW:] DL/T 5429-2009、远景展望、水平年、10\textasciitilde{}15年、电力系统设计
  \item[标准:] DL/T 5429-2009
  \item[条款溯源:] DL/T 5429-2009 第3.0.7条
  \item[主题:] 通用
  \item[评分方式:] auto\_keyword\_match
  \item[Rubric条款:] DL/T 5429-2009 第3.0.7条
  \item[Rubric判据:] 答案应为: 10\textasciitilde{}15年

\end{description}
\vspace{4pt}
\hrule
\vspace{6pt}

\noindent\textbf{L1-053} \hfill \textbf{参数检索} $|$ \textit{数值阈值与限值}
\begin{description}[leftmargin=1.5em,style=nextline,font=\normalfont]
  \item[Q:] 220kV变电站的220kV配电装置，出线回路数在多少回及以下时可采用其他简单主接线？
  \item[A:] 4回及以下
  \item[KW:] 220kV变电站、220kV配电装置、出线回路数、4回、简单主接线
  \item[标准:] DL/T 5218-2012
  \item[条款溯源:] DL/T 5218-2012 第5.1.6条
  \item[主题:] 通用
  \item[评分方式:] auto\_keyword\_match
  \item[Rubric条款:] DL/T 5218-2012 第5.1.6条
  \item[Rubric判据:] 答案应为: 4回及以下

\end{description}
\vspace{4pt}
\hrule
\vspace{6pt}

\noindent\textbf{L1-014} \hfill \textbf{参数检索} $|$ \textit{无功补偿参数}
\begin{description}[leftmargin=1.5em,style=nextline,font=\normalfont]
  \item[Q:] 根据GB 38755-2019，电网的无功补偿应遵循什么具体比例原则？
  \item[A:] 分层分区和就地平衡
  \item[KW:] 无功补偿、分层分区、就地平衡、原则、GB 38755-2019
  \item[标准:] GB 38755-2019 §3.4.2
  \item[条款溯源:] GB 38755-2019 §3.4.2
  \item[主题:] 电压稳定
  \item[评分方式:] auto\_keyword\_match
  \item[Rubric条款:] GB 38755-2019 §3.4.2
  \item[Rubric判据:] 答案应为: 分层分区和就地平衡

\end{description}
\vspace{4pt}
\hrule
\vspace{6pt}

\noindent\textbf{L1-005} \hfill \textbf{参数检索} $|$ \textit{保护配置要求}
\begin{description}[leftmargin=1.5em,style=nextline,font=\normalfont]
  \item[Q:] 根据GB 38755-2019，电力系统发生稳定破坏时，必须采取预定的措施以防止事故范围扩大的时间要求是多少？
  \item[A:] 无具体数值参数
  \item[KW:] 稳定破坏、预定措施、事故范围、扩大、减少损失
  \item[标准:] GB 38755-2019 §3.1.5
  \item[条款溯源:] GB 38755-2019 4.2.4
  \item[主题:] 通用
  \item[评分方式:] auto\_keyword\_match
  \item[Rubric条款:] GB 38755-2019 4.2.4
  \item[Rubric判据:] 答案应为: 无具体数值参数

\end{description}
\vspace{4pt}
\hrule
\vspace{6pt}

\noindent\textbf{L1-049} \hfill \textbf{参数检索} $|$ \textit{设备参数与容量}
\begin{description}[leftmargin=1.5em,style=nextline,font=\normalfont]
  \item[Q:] DL/T 5218-2012中，每台站用变压器容量按什么选择？
  \item[A:] 按全站计算负荷选择
  \item[KW:] DL/T 5218-2012、站用变压器、容量、全站计算负荷、选择
  \item[标准:] DL/T 5218-2012 §3.0.9
  \item[条款溯源:] DL/T 5218-2012 §5.6.2
  \item[主题:] 变电站设计
  \item[评分方式:] auto\_keyword\_match
  \item[Rubric条款:] DL/T 5218-2012 §5.6.2
  \item[Rubric判据:] 答案应为: 按全站计算负荷选择

\end{description}
\vspace{4pt}
\hrule
\vspace{6pt}

\noindent\textbf{L1-069} \hfill \textbf{参数检索} $|$ \textit{数值阈值与限值}
\begin{description}[leftmargin=1.5em,style=nextline,font=\normalfont]
  \item[Q:] 根据DL/T 5429-2009，系统总备用容量可按系统最大发电负荷的多少比例考虑？
  \item[A:] 15\%\textasciitilde{}20\%
  \item[KW:] 系统总备用容量、最大发电负荷、15\%、20\%、比例
  \item[标准:] DL/T 5429-2009
  \item[条款溯源:] DL/T 5429-2009 第5.2.3条
  \item[主题:] 通用
  \item[评分方式:] auto\_keyword\_match
  \item[Rubric条款:] DL/T 5429-2009 第5.2.3条
  \item[Rubric判据:] 答案应为: 15\%\textasciitilde{}20\%

\end{description}
\vspace{4pt}
\hrule
\vspace{6pt}

\noindent\textbf{L1-072} \hfill \textbf{参数检索} $|$ \textit{数值阈值与限值}
\begin{description}[leftmargin=1.5em,style=nextline,font=\normalfont]
  \item[Q:] 根据DL/T 5429-2009，电力系统设计的远景展望年限具体是多少年？
  \item[A:] 10～15年
  \item[KW:] DL/T 5429-2009、电力系统设计、远景展望、年限、10～15年
  \item[标准:] DL/T 5429-2009
  \item[条款溯源:] DL/T 5429-2009 第3.0.7条
  \item[主题:] 通用
  \item[评分方式:] auto\_keyword\_match
  \item[Rubric条款:] DL/T 5429-2009 第3.0.7条
  \item[Rubric判据:] 答案应为: 10～15年

\end{description}
\vspace{4pt}
\hrule
\vspace{6pt}

\section*{L2 — 推理型（需条件判断或跨参数推理）}

\noindent\textbf{L2-043} \hfill \textbf{审查判断} $|$ \textit{稳定校核}
\begin{description}[leftmargin=1.5em,style=nextline,font=\normalfont]
  \item[Q:] 【场景】某500kV输电线路同塔双回架设，根据DL/T 5218-2012《220kV\textasciitilde{}750kV变电站设计规程》，需考虑双回线同时停运的N-2故障。该线路输送功率占受端系统总负荷的40\%，且无其他备用联络线。

【问题】请分析双回线同停对受端系统电压和频率的影响。若系统频率下降至49.2Hz，根据规程要求，应启动多少容量的低频减载？请计算并说明依据。
  \item[A:] 双回线同停导致受端系统损失40\%电源，频率和电压大幅下降。根据DL/T 5218-2012，频率降至49.2Hz时，应切除不少于总负荷的30\%以恢复频率。假设系统总负荷为1000MW，则需切除300MW。同时需考虑电压崩溃风险，可能需联切部分负荷。
  \item[KW:] N-2故障、双回线同停、低频减载、频率稳定、电压崩溃
  \item[标准:] DL/T 5218-2012 §7.3.1
  \item[条款溯源:] DL/T 5218-2012 §7.3.1
  \item[主题:] 电网结构
  \item[评分方式:] manual\_review
  \item[Rubric条款:] DL/T 5218-2012
  \item[Rubric判据:] 2Hz时，应切除不少于总负荷的30\%以恢复频率、假设系统总负荷为1000MW，则需切除300MW、同时需考虑电压崩溃风险，可能需联切部分负荷

\end{description}
\vspace{4pt}
\hrule
\vspace{6pt}

\noindent\textbf{L2-044} \hfill \textbf{审查判断} $|$ \textit{稳定校核}
\begin{description}[leftmargin=1.5em,style=nextline,font=\normalfont]
  \item[Q:] 【场景】某220kV变电站有两台主变压器，容量均为180MVA，短路阻抗百分比为14\%。正常运行时，每台主变负载率为65\%。现计划扩建，新增一台同容量主变，三台并列运行。根据GB 38755-2019《电力系统安全稳定导则》，需进行N-1校验。

【问题】若其中一台主变因故障退出运行，剩余两台主变能否满足N-1后不过载要求？请计算剩余两台主变的负载率，并判断是否超过1.3倍额定容量限额。
  \item[A:] 正常三台总负荷为3×180×0.65=351MVA。一台退出后，剩余两台承担351MVA，每台负载率=351/(2×180)=97.5\%。根据GB 38755-2019，N-1后允许短时过载不超过1.3倍额定容量，即1.3×180=234MVA，实际每台180×0.975=175.5MVA，未超过限额，满足要求。
  \item[KW:] 负载率计算、N-1校验、过载限额、主变压器、短时过载
  \item[标准:] GB 38755-2019 §4.2.2
  \item[条款溯源:] GB 38755-2019 第2.3条
  \item[主题:] N-1原则
  \item[评分方式:] manual\_review
  \item[Rubric条款:] GB 38755-2019
  \item[Rubric判据:] 根据GB 38755-2019，N-1后允许短时过载不超过1

\end{description}
\vspace{4pt}
\hrule
\vspace{6pt}

\noindent\textbf{L2-007} \hfill \textbf{审查判断} $|$ \textit{无功补偿}
\begin{description}[leftmargin=1.5em,style=nextline,font=\normalfont]
  \item[Q:] 【场景】某500kV超高压输电工程中，一条长度为180km的架空线路采用4×LGJ-400/35钢芯铝绞线，单位长度充电功率为0.833Mvar/km，总充电功率为150Mvar。该线路通过一台容量为750MVA的500/220kV自耦变压器向220kV电网供电，变压器短路阻抗为14\%，220kV侧最大负荷为500MW，功率因数为0.95。该500kV线路末端装设有两组容量分别为90Mvar和60Mvar的并联电抗器，总补偿容量为150Mvar。根据GB 38755-2019第3.4.2条关于无功补偿的原则，判断该线路的充电功率是否需要补偿？并说明理由。

【问题】某500kV电网中，一条500kV线路的充电功率为150Mvar，该线路通过一台500/220kV变压器向220kV电网供电。根据GB 38755-2019第3.4.2条关于无功补偿的原则，判断该线路的充电功率是否需要补偿？并说明理由。
  \item[A:] 需要补偿。根据GB 38755-2019 §3.4.2，330kV及以上等级架空线路的充电功率应基本予以补偿。该线路为500kV架空线路，属于330kV及以上等级，其充电功率150Mvar应基本予以补偿，以避免无功功率长距离传输，减少线损，控制电压水平。补偿方式可采用并联电抗器。
  \item[KW:] 充电功率、补偿、330kV及以上、并联电抗器
  \item[标准:] GB 38755-2019 §3.4.2
  \item[条款溯源:] GB 38755-2019 第3.4.2条
  \item[主题:] 电压稳定
  \item[评分方式:] manual\_review
  \item[Rubric条款:] GB 38755-2019 §3.4.2、§3.4.2
  \item[Rubric判据:] 2，330kV及以上等级架空线路的充电功率应基本予以补偿、该线路为500kV架空线路，属于330kV及以上等级，其充电功率150Mvar应基本予以补偿，以避免无功功率长距离传输，减少线损，控制电压水平
  \item[LJ判定:] 需重写 (弱=0.60, 强=0.90, 差距=0.30)
  \item[LJ反馈:] 题目推理深度不足：弱模型得分=0.60≥0.4。请增加推理复杂度，加入更多条件判断和参数交互。
\end{description}
\vspace{4pt}
\hrule
\vspace{6pt}

\noindent\textbf{L2-021} \hfill \textbf{审查判断} $|$ \textit{电网结构}
\begin{description}[leftmargin=1.5em,style=nextline,font=\normalfont]
  \item[Q:] 【场景】某沿海规划中的受端电网，其总负荷为20000MW，外部电源通过三组独立的500kV送电回路接入该受端系统。每组送电回路采用双回500kV线路，单回线路的送电能力按500kV电压等级典型参数取1200MW，双回合计可达2400MW。受端系统内配置有3台容量为1000MVA的主变压器，短路电流水平按500kV母线短路容量计算约为40kA。该受端系统缺乏动态无功支撑，且直流落点集中，仅配置了总容量为400Mvar的容性无功补偿装置，远低于标准要求的无功补偿容量。根据GB 38755-2019 §3.2.2.2的要求，外部电源需经相对独立的送电回路接入受端系统，避免送电回路落点和输电走廊过于集中。请结合上述工程参数和标准要求，判断每组送电回路的最大输送功率应控制在什么范围？并说明理由。

【问题】原问题文本：某规划中的受端系统总负荷为20000MW，外部电源通过三组独立的送电回路接入该系统。根据GB 38755-2019，每组送电回路的最大输送功率占受端系统总负荷的比例不应过大，需结合具体条件决定。若该受端系统缺乏动态无功支撑，且直流落点集中，请判断每组送电回路的最大输送功率应控制在什么范围？并说明理由。
  \item[A:] 每组送电回路的最大输送功率应控制在2000MW至4000MW之间（即总负荷的10\%-20\%）。理由：根据GB 38755-2019 §3.2.2.2，外部电源需经相对独立的送电回路接入受端系统，避免送电回路落点和输电走廊过于集中。当受端系统缺乏动态无功支撑且直流落点集中时，系统承受大功率冲击的能力较弱，若单组送电回路功率过大，一旦故障断开可能导致受端系统频率和电压大幅波动，甚至引发稳定破坏。因此，应取较低的比例（如10\%-15\%），确保任一送电回路断开后，剩余电源和本地机组能够维持系统稳定，且其他元件不过负荷。
  \item[KW:] 受端系统、送电回路、最大输送功率、动态无功支撑、稳定破坏
  \item[标准:] GB 38755-2019 §3.2.2.2
  \item[条款溯源:] GB 38755-2019 第3.2.2.2条
  \item[主题:] 电网结构
  \item[评分方式:] manual\_review
  \item[Rubric条款:] GB 38755-2019 §3.2.2.2、§3.2.2.2
  \item[Rubric判据:] 每组送电回路的最大输送功率应控制在2000MW至4000MW之间（即总负荷的10\%-20\%）、2，外部电源需经相对独立的送电回路接入受端系统，避免送电回路落点和输电走廊过于集中、因此，应取较低的比例（如10\%-15\%），确保任一送电回路断开后，剩余电源和本地机组能够维持系统稳定，且其他元件不过负荷
  \item[LJ判定:] 需重写 (弱=0.68, 强=0.50, 差距=-0.18)
  \item[LJ反馈:] 题目推理深度不足：弱模型得分=0.68≥0.4。请增加推理复杂度，加入更多条件判断和参数交互。
\end{description}
\vspace{4pt}
\hrule
\vspace{6pt}

\noindent\textbf{L2-018} \hfill \textbf{审查判断} $|$ \textit{短路电流}
\begin{description}[leftmargin=1.5em,style=nextline,font=\normalfont]
  \item[Q:] 【场景】某330kV枢纽变电站位于区域负荷中心，通过4回300km输电线路与主力电厂相连，主变容量为3×240MVA，无功补偿容量为200Mvar。该变电站母线短路电流水平经实测为50kA，而原设计阶段基于电网规划确定的短路电流耐受能力为40kA。当前该站供电区域负荷持续增长，系统联络线紧密，短路电流水平呈上升趋势。根据GB 38755-2019《电力系统安全稳定导则》第3.2.4.3条关于分区电网应有效限制短路电流的要求，以及第3.1.2条e)款关于合理控制系统短路电流的规定，需对该变电站短路电流超标问题进行评估。

【问题】请分析该变电站是否满足标准要求？若不满足，标准建议采取哪些措施？
  \item[A:] 根据GB 38755-2019 §3.2.4.3，分区电网应尽可能简化，以有效限制短路电流和简化继电保护的配置。本例中，实际短路电流50kA超过设计耐受能力40kA，不满足要求。标准§3.1.2 e)要求合理控制系统短路电流。应采取的措施包括：调整电网运行方式（如断开部分联络线）、加装串联电抗器、采用高阻抗变压器、或分层分区运行以降低短路电流水平。
  \item[KW:] 限制短路电流、50kA、40kA、串联电抗器、分层分区
  \item[标准:] GB 38755-2019 §3.2.4.3 \& §3.1.2 e)
  \item[条款溯源:] GB 38755-2019 第3.2.4.3条, GB 38755-2019 第3.1.2条e)款
  \item[主题:] 短路电流
  \item[评分方式:] manual\_review
  \item[Rubric条款:] GB 38755-2019 §3.2.4.3、§3.2.4.3、§3.1.2
  \item[Rubric判据:] 3，分区电网应尽可能简化，以有效限制短路电流和简化继电保护的配置、本例中，实际短路电流50kA超过设计耐受能力40kA，不满足要求、2 e)要求合理控制系统短路电流、应采取的措施包括：调整电网运行方式（如断开部分联络线）、加装串联电抗器、采用高阻抗变压器、或分层分区运行以降低短路电流水平
  \item[LJ判定:] 需重写 (弱=0.54, 强=0.77, 差距=0.23)
  \item[LJ反馈:] 题目推理深度不足：弱模型得分=0.54≥0.4。请增加推理复杂度，加入更多条件判断和参数交互。
\end{description}
\vspace{4pt}
\hrule
\vspace{6pt}

\noindent\textbf{L2-072} \hfill \textbf{审查判断} $|$ \textit{过电压保护}
\begin{description}[leftmargin=1.5em,style=nextline,font=\normalfont]
  \item[Q:] 【场景】某500kV变电站过电压保护设计中，避雷器配置不足导致雷击跳闸。设计人员需依据DL/T 5218-2012第9章（过电压保护）重新计算保护范围。

【问题】在500kV变电站中，如何根据DL/T 5218-2012第9章确定避雷器的安装位置和保护范围？请结合雷击跳闸案例说明。
  \item[A:] 根据DL/T 5218-2012第9章，避雷器应安装在变压器、母线及线路入口处，保护范围通过雷电侵入波计算确定。典型配置为：主变压器侧装设一组避雷器，母线装设一组，线路入口装设一组。保护距离应小于避雷器额定电压下的最大允许距离，否则需增设避雷器。
  \item[KW:] 场景、变电站过、电压保护、设计中、避雷器配
  \item[标准:] DL/T 5218-2012
  \item[条款溯源:] DL/T 5218-2012 第9章
  \item[主题:] 过电压保护
  \item[评分方式:] manual\_review
  \item[Rubric条款:] DL/T 5218-2012
  \item[Rubric判据:] 根据DL/T 5218-2012第9章，避雷器应安装在变压器、母线及线路入口处，保护范围通过雷电侵入波计算确定、保护距离应小于避雷器额定电压下的最大允许距离，否则需增设避雷器

\end{description}
\vspace{4pt}
\hrule
\vspace{6pt}

\noindent\textbf{L2-090} \hfill \textbf{审查判断} $|$ \textit{保护配置}
\begin{description}[leftmargin=1.5em,style=nextline,font=\normalfont]
  \item[Q:] 【场景】某500kV变电站采用无人值班运行管理模式，根据DL/T 5218-2012第8章电气二次要求，需配置计算机监控系统。设计人员需确定监控系统的功能范围，包括遥测、遥信、遥控、遥调。请分析在无人值班模式下，监控系统应具备哪些关键功能？并说明对继电保护信息管理的要求。

【问题】在上述场景下，请分析监控系统需具备哪些关键功能，并判断继电保护信息管理应满足哪些具体要求？
  \item[A:] 根据DL/T 5218-2012第8章，无人值班变电站计算机监控系统应具备以下关键功能：1）遥测：采集母线电压、线路电流、有功无功功率、变压器油温等模拟量；2）遥信：采集断路器、隔离开关位置信号，保护动作信号，设备状态信号；3）遥控：远程控制断路器分合、隔离开关操作、电容器组投切；4）遥调：调节变压器分接头、无功补偿装置出力。对继电保护信息管理的要求：监控系统应能接收保护装置的故障录波数据、动作事件记录、定值信息，并具备远方修改保护定值功能（需经安全认证），同时应实现保护装置与监控系统的时间同步，确保故障分析准确性。
  \item[KW:] 场景、变电站采、用无人值、班运行管、理模式
  \item[标准:] DL/T 5218-2012
  \item[条款溯源:] DL/T 5218-2012 第8章
  \item[主题:] 电气主接线
  \item[评分方式:] manual\_review
  \item[Rubric条款:] DL/T 5218-2012
  \item[Rubric判据:] 根据DL/T 5218-2012第8章，无人值班变电站计算机监控系统应具备以下关键功能：1）遥测：采集母线电压、线路电流、有功无功功率、变压器油温等模拟量、对继电保护信息管理的要求：监控系统应能接收保护装置的故障录波数据、动作事件记录、定值信息，并具备远方修改保护定值功能（需经安全认证），同时应实现保护装置与监控系统的时间同步，确保故障分析准确性

\end{description}
\vspace{4pt}
\hrule
\vspace{6pt}

\noindent\textbf{L2-032} \hfill \textbf{审查判断} $|$ \textit{电网结构}
\begin{description}[leftmargin=1.5em,style=nextline,font=\normalfont]
  \item[Q:] 【场景】某省级电网规划建设一座500kV枢纽变电站，该站同时作为两个直流输电工程的落点换流站，每个直流工程额定输送容量均为3000MW。该受端系统最高一级电压为500kV，区域内已有多座大型火电厂和抽水蓄能电站，系统短路电流水平在40kA至50kA之间。该枢纽变电站规划安装两台容量为1000MVA的主变压器，并配置±300Mvar的动态无功补偿装置。根据GB 38755-2019中3.2.1条的要求，请分析该枢纽变电站的规划方案应满足哪些条件？若该站仅通过单回500kV线路与主网相连，该单回线路的送电能力为2000MW，是否合理？为什么？
  \item[A:] 应满足：1)枢纽变电站的规模和换流站的容量应与受端系统相适应；2)直流落点与负荷集中地区应合理配置动态无功调节设备；3)应优化直流落点，完善近区网架，提高系统对直流的支撑能力。单回500kV线路不合理，因为枢纽变电站应具有较坚强的网络联系，单回线路不满足N-1原则，且无法满足多馈入直流对系统支撑能力的要求。
  \item[KW:] 枢纽变电站、受端系统、动态无功调节、多馈入直流、N-1原则
  \item[标准:] GB 38755-2019 §3.2.1
  \item[条款溯源:] GB 38755-2019 第3.2.1条
  \item[主题:] 直流输电
  \item[评分方式:] manual\_review
  \item[Rubric条款:] GB 38755-2019 §3.2.1
  \item[Rubric判据:] 应满足：1)枢纽变电站的规模和换流站的容量应与受端系统相适应、2)直流落点与负荷集中地区应合理配置动态无功调节设备、3)应优化直流落点，完善近区网架，提高系统对直流的支撑能力、单回500kV线路不合理，因为枢纽变电站应具有较坚强的网络联系，单回线路不满足N-1原则，且无法满足多馈入直流对系统支撑能力的要求

\end{description}
\vspace{4pt}
\hrule
\vspace{6pt}

\noindent\textbf{L2-001} \hfill \textbf{审查判断} $|$ \textit{电网结构}
\begin{description}[leftmargin=1.5em,style=nextline,font=\normalfont]
  \item[Q:] 【场景】某沿海经济新区规划建设500kV受端电网，该区域负荷集中，预计2025年最大负荷达8000MW，主要由电子信息、高端制造等产业构成，负荷特性对电压稳定性要求较高。外部电源通过三回独立的500kV交流线路向该受端电网送电，每回线路采用4×LGJ-400/50钢芯铝绞线，线路长度约180km，根据参数约束表，500kV线路单回最大送电能力可达1200MW，三回总送电能力为3600MW。受端电网内规划新建两座500kV变电站，主变容量均为1000MVA，短路电流水平按40kA设计，并配置±200Mvar的STATCOM无功补偿装置以支撑电压。该区域负荷同时由本地一座600MW火电机组和一座400MW燃气机组支撑，但外部电源仍占主导地位。根据GB 38755-2019第3.2.2.2条关于外部电源送电回路的规定，判断该送电方案是否满足安全稳定要求？并说明理由。
  \item[A:] 不满足要求。根据GB 38755-2019 §3.2.2.2，每一组送电回路的最大输送功率所占受端系统总负荷的比例不应过大。三回线路总送电功率为3600MW，占受端总负荷8000MW的比例为45\%。虽然标准未给出具体比例数值，但要求结合受端系统具体条件决定，通常认为单组或多组送电回路总容量不宜超过受端负荷的30\%-40\%。45\%的比例过高，一旦多回线路同时故障，将导致受端系统出现严重功率缺额，威胁安全稳定。
  \item[KW:] 送电回路、最大输送功率、受端系统总负荷、比例不应过大
  \item[标准:] GB 38755-2019 §3.2.2.2
  \item[条款溯源:] GB 38755-2019 §3.2.2.2
  \item[主题:] 电网结构
  \item[评分方式:] manual\_review
  \item[Rubric条款:] GB 38755-2019 §3.2.2.2、§3.2.2.2
  \item[Rubric判据:] 2，每一组送电回路的最大输送功率所占受端系统总负荷的比例不应过大、虽然标准未给出具体比例数值，但要求结合受端系统具体条件决定，通常认为单组或多组送电回路总容量不宜超过受端负荷的30\%-40\%
  \item[LJ判定:] 通过 (弱=0.10, 强=0.40, 差距=0.30)
\end{description}
\vspace{4pt}
\hrule
\vspace{6pt}

\noindent\textbf{L2-056} \hfill \textbf{审查判断} $|$ \textit{稳定校核}
\begin{description}[leftmargin=1.5em,style=nextline,font=\normalfont]
  \item[Q:] 【场景】某沿海火电厂装机容量为900MW，以两回220kV交流线路接入500kV枢纽变电站。220kV线路采用2×LGJ-400/50钢芯铝绞线，每回线路额定送电能力为450MW（基于环境温度40℃、导线允许温度80℃条件下计算）。电厂220kV母线短路电流水平为38kA，主变容量为3×360MVA。根据DL/T 5429-2009规程要求，对该电厂送出线路进行安全稳定校核。正常方式下，两回线路并列运行，总输电容量为800MW（含厂用电及损耗）。现假设一回线路发生三相短路故障，保护动作切除该线路，事故后剩余一回线路的静稳定能力经计算为700MW（考虑线路充电功率及无功补偿装置投切后的极限输送能力）。电厂侧配置有快速切机装置，可在0.2秒内切除部分机组出力。

【问题】请判断该电厂送出线路是否满足安全稳定要求，并说明是否需要采取切机措施。
  \item[A:] 规程要求电厂送出线路两回及以上时，任一回线路事故停运后，若事故后静稳定能力小于正常输电容量，按事故后静稳定能力输电。本例事故后静稳定能力700MW小于正常输电容量800MW，因此应按700MW输电，即需将出力降至700MW以下，需采取切机或快速降出力措施。
  \item[KW:] 事故后静稳定能力、切机措施、送出线路、火电厂、DL/T 5429-2009 §6.3.3
  \item[标准:] DL/T 5429-2009 §6.3.3
  \item[条款溯源:] DL/T 5429-2009 §6.3.3
  \item[主题:] 安全自动装置
  \item[评分方式:] manual\_review
  \item[Rubric条款:] DL/T 5429-2009 §6.3.3
  \item[Rubric判据:] 规程要求电厂送出线路两回及以上时，任一回线路事故停运后，若事故后静稳定能力小于正常输电容量，按事故后静稳定能力输电、本例事故后静稳定能力700MW小于正常输电容量800MW，因此应按700MW输电，即需将出力降至700MW以下，需采取切机或快速降出力措施

\end{description}
\vspace{4pt}
\hrule
\vspace{6pt}

\section*{L3 — 综合型（需引用两份以上规范综合判断，五段式结构）}

\noindent\textbf{L3-100} \hfill \textbf{综合评估} $|$ \textit{综合决策-继电保护}
\begin{description}[leftmargin=1.5em,style=nextline,font=\normalfont]
  \item[Q:] 某大型燃煤电厂现有装机容量为1200MW（4×300MW机组），通过一座220kV升压站以2回220kV输电线路接入系统，线路长度分别为45km和52km，导线型号为2×LGJ-400/50，正常输送容量约500MW，220kV母线短路电流水平为28kA。该升压站运行已超过15年，现有设备为常规电磁式电流/电压互感器、集中式监控系统及基于硬接线的继电保护装置，保护动作时间约40ms，故障录波精度为1ms。计划对该220kV升压站进行智能化改造，同时需满足GB 38755-2019《电力系统安全稳定导则》对系统稳定性的要求（如故障清除时间要求不大于100ms、暂态稳定裕度不低于15\%），以及DL/T 5218-2012《220kV\textasciitilde{}750kV变电站设计技术规程》对变电站设计的技术规定（包括防误操作、防越级跳闸、二次回路抗干扰能力及信息共享要求）。改造方案中，A方案为采用常规保护与测控装置，保留原有集中式监控系统，仅更换部分老旧电磁式互感器为同类设备，保护动作时间可优化至30ms，但无法实现自适应整定和网络化信息共享；B方案为全面采用数字化保护、测控及智能终端，构建全站IEC 61850通信网络，配置电子式互感器（采样精度0.2级，保护动作时间可降至15ms），并部署在线监测与智能告警系统，实现设备状态实时感知及故障快速定位。请基于上述两个标准，从系统安全稳定、设计合规性、运行可靠性及运维效率等维度，对A、B方案进行综合评估，并给出推荐方案及控制优先级链条。
  \item[A:] 现状诊断：该电厂220kV升压站运行年限较长，现有设备为常规电磁式互感器与集中式监控系统，保护动作时间及故障录波精度已无法满足GB 38755-2019对系统故障清除时间及动态稳定控制的要求；同时，DL/T 5218-2012要求变电站设计应具备完善的防误操作、防越级跳闸及抗干扰能力，现有设计在二次回路抗干扰及信息共享方面存在短板，亟需通过改造提升整体安全稳定水平。

多维冲突分析：A方案（常规改造）与B方案（数字化改造）存在以下冲突：一是安全稳定维度，GB 38755-2019强调故障快速切除与系统动态稳定，A方案依赖传统硬接线，保护动作时间较长且缺乏自适应整定能力，难以满足暂态稳定要求；B方案通过数字采样与GOOSE跳闸可显著缩短动作时间，但需解决网络延时及同步问题。二是设计合规维度，DL/T 5218-2012要求二次设备抗干扰能力及冗余配置，A方案沿用常规设计，抗干扰措施成熟但信息共享差；B方案需额外满足智能电子设备（IED）的电磁兼容及网络冗余要求，设计复杂度高。三是运维效率维度，A方案维护依赖人工巡检，故障定位慢；B方案支持在线监测与自诊断，但初期投资大且对运维人员技能要求高。

多方案比选：A方案优势在于技术成熟、投资低、改造周期短，且与现有运行规程兼容性好；劣势是难以实现系统级稳定控制优化，且无法满足未来智能化调度需求。B方案优势在于符合GB 38755-2019对快速故障隔离及广域保护的要求，同时满足DL/T 5218-2012对数字化设计的前瞻性要求，可提升运行可靠性并降低全生命周期运维成本；劣势是初期投资高、需解决网络风暴及时钟同步等风险，且改造期间需严格管控停电范围。综合评估，B方案在安全稳定与设计合规性上更优，但需通过分阶段实施及冗余配置来规避风险。

折中综合方案：推荐采用B方案，但采取分阶段过渡策略：第一阶段保留现有母线保护及主变压器保护的硬接线回路作为后备，新增数字化保护与测控系统，采用双网冗余及IEEE 1588精确对时；第二阶段逐步替换老旧互感器为电子式互感器，并接入在线监测系统；第三阶段实现全站智能化运维。同时，在设计中严格按DL/T 5218-2012要求配置防跳、防误闭锁及抗干扰措施，并依据GB 38755-2019进行系统稳定校核，确保改造后故障清除时间不大于100ms。

控制优先级链条：第一优先级：确保故障快速切除与系统稳定（GB 38755-2019第3.2条）→第二优先级：二次设备抗干扰与冗余设计（DL/T 5218-2012第5.4条）→第三优先级：数字化通信网络可靠性（IEC 61850一致性测试）→第四优先级：在线监测与智能告警功能投运→第五优先级：运维规程修订与人员培训。
  \item[KW:] GB 38755-2019、DL/T 5218-2012、安全稳定、数字化改造、控制优先级
  \item[标准:] GB 38755-2019, DL/T 5218-2012
  \item[条款溯源:] GB 38755-2019 第4.2.4条
  \item[主题:] 通用
  \item[评分方式:] manual\_review
  \item[Rubric条款:] GB 38755-2019、DL/T 5218-2012、IEEE 1588、IEC 61850
  \item[Rubric判据:] 现状诊断：该电厂220kV升压站运行年限较长，现有设备为常规电磁式互感器与集中式监控系统，保护动作时间及故障录波精度已无法满足GB 38755-2019对系统故障清除时间及动态稳定控制的要求、同时，DL/T 5218-2012要求变电站设计应具备完善的防误操作、防越级跳闸及抗干扰能力，现有设计在二次回路抗干扰及信息共享方面存在短板，亟需通过改造提升整体安全稳定水平、多维冲突分析：A方案（常规改造）与B方案（数字化改造）存在以下冲突：一是安全稳定维度，GB 38755-2019强调故障快速切除与系统动态稳定，A方案依赖传统硬接线，保护动作时间较长且缺乏自适应整定能力，难以满足暂态稳定要求、B方案通过数字采样与GOOSE跳闸可显著缩短动作时间，但需解决网络延时及同步问题、二是设计合规维度，DL/T 5218-2012要求二次设备抗干扰能力及冗余配置，A方案沿用常规设计，抗干扰措施成熟但信息共享差
  \item[LJ判定:] 通过 (弱=0.80, 强=0.80, 差距=0.00)
\end{description}
\vspace{4pt}
\hrule
\vspace{6pt}

\noindent\textbf{L3-012} \hfill \textbf{综合评估} $|$ \textit{方案比选-接入系统}
\begin{description}[leftmargin=1.5em,style=nextline,font=\normalfont]
  \item[Q:] 【场景】某沿海省份电网规划建设一座新的核电站（2×1000MW），并配套接入500kV主网。核电站附近有大型风电场（800MW）和一座抽水蓄能电站（4×300MW）。现有两个接入方案：

【问题】
方案A：核电站通过双回500kV线路直接接入枢纽变电站，线路长度80km，不设开关站。核电机组功率因数按0.85（滞相）设计，不参与调峰。
方案B：核电站通过单回500kV线路接入新建的500kV开关站，再通过双回线路接入主网，线路总长100km。核电机组功率因数按0.95（滞相）设计，并预留10\%的出力用于一次调频。

请依据GB 38755-2019和DL/T 5429-2009，对上述两个方案进行综合评估，并给出推荐方案及控制优先级链条。
  \item[A:] 现状诊断：该区域核电、风电、抽蓄并存，电源结构复杂。根据DL/T 5429-2009第5.1.3条，核电站接入系统需满足更高的安全稳定标准。GB 38755-2019第3.3条要求电源结构应合理配置，确保系统频率和电压稳定。当前问题在于核电站的接入方式、无功配置及调峰调频能力对系统安全稳定的影响。

多维冲突分析：
1. 安全稳定与接入可靠性的冲突：方案A双回直连枢纽站，可靠性高，但根据GB 38755-2019第4.2.2条，若双回线同时故障（如台风导致），将损失2000MW核电，需校核系统频率稳定性。方案B通过开关站增加一级冗余，但单回线至开关站故障时，仍可能损失全部核电。
2. 无功平衡与电压稳定的冲突：方案A核电机组功率因数0.85，可提供大量无功支撑，符合GB 38755-2019第3.4条无功补偿要求，但滞相运行增加发电机损耗。方案B功率因数0.95，无功支撑能力弱，需额外配置高压并联电抗器或STATCOM，增加投资。
3. 调峰调频与经济性的冲突：方案A核电机组不参与调峰，需依赖抽水蓄能电站承担全部调峰任务（1200MW），根据DL/T 5429-2009第5.3.1条，需校核抽蓄容量是否足够。方案B预留10\%出力调频，但核电机组频繁调节可能影响堆芯安全，且经济性下降。

多方案比选：
- 方案A优势：接入简单，投资低；核电无功支撑强，可减少系统无功补偿投资。劣势：核电机组不参与调频，系统频率稳定性依赖抽蓄快速响应（GB 38755-2019第5.7节）；双回线同时故障风险高。
- 方案B优势：通过开关站增加运行灵活性；核电机组参与一次调频，提升频率稳定性。劣势：投资高；核电无功支撑弱，需额外配置无功补偿（如2组120Mvar并联电抗器）；核电机组调频可能影响寿命。

折中综合方案：推荐采用方案A的接入方式，但附加以下条件：
1. 根据GB 38755-2019第4.2.2条，配置抽水蓄能电站的快速切泵功能，在核电跳闸后0.5s内切除300MW抽水负荷，维持功率平衡。
2. 根据DL/T 5429-2009第5.1.3条，要求核电机组具备0.90（滞相）的功率因数运行能力，并配置自动电压调节器（AVR），确保电压稳定。
3. 在核电站500kV出线侧配置一组180Mvar的高压并联电抗器，用于轻载时吸收无功，防止电压升高（GB 38755-2019第3.4条）。

控制优先级链条：
第一优先级：抽水蓄能电站快速切泵，防止频率崩溃（GB 38755-2019第5.7节）。
第二优先级：核电机组AVR动作，维持电压稳定（GB 38755-2019第3.4条）。
第三优先级：高压并联电抗器自动投切，控制无功平衡（DL/T 5429-2009第6.1.6条）。
第四优先级：风电场的无功补偿装置（SVG）动态响应，提供电压支撑。
  \item[KW:] GB 38755-2019第4.2.2条、DL/T 5429-2009第5.1.3条、无功分层分区平衡、一次调频、抽水蓄能切泵
  \item[标准:] GB 38755-2019, DL/T 5429-2009
  \item[条款溯源:] GB 38755-2019 第4.2.2条
  \item[主题:] 电压稳定
  \item[评分方式:] manual\_review
  \item[Rubric条款:] DL/T 5429-2009、GB 38755-2019、GB 38755-2019 第3.3条、GB 38755-2019 第4.2.2条、GB 38755-2019 第3.4条、GB 38755-2019 第5.7节、DL/T 5429-2009 第5.1.3条、DL/T 5429-2009 第5.3.1条
  \item[Rubric判据:] 3条，核电站接入系统需满足更高的安全稳定标准、3条要求电源结构应合理配置，确保系统频率和电压稳定、2条，若双回线同时故障（如台风导致），将损失2000MW核电，需校核系统频率稳定性、85，可提供大量无功支撑，符合GB 38755-2019第3、4条无功补偿要求，但滞相运行增加发电机损耗
  \item[LJ判定:] 需重写 (弱=0.40, 强=0.40, 差距=0.00)
  \item[LJ反馈:] L3区分度不足：差距=0.00。请增强方案间的技术冲突或增加工程决策的难度。
\end{description}
\vspace{4pt}
\hrule
\vspace{6pt}

\noindent\textbf{L3-077} \hfill \textbf{综合评估} $|$ \textit{方案比选-变电站设计}
\begin{description}[leftmargin=1.5em,style=nextline,font=\normalfont]
  \item[Q:] 【场景】某500kV枢纽变电站拟进行电气主接线及配电装置设计，站址位于沿海高盐雾地区，地震设防烈度为8度。该变电站规划安装2台1200MVA主变压器，远期规模为4台1200MVA主变压器；500kV出线规划8回，本期建设4回，每回线路输送容量约1500MW，最大短路电流按50kA考虑；220kV出线规划16回，本期建设8回，每回线路输送容量约400MW，最大短路电流按40kA考虑。站址所在区域年平均盐雾密度高达0.5mg/cm²，属于严重污秽等级，且8度地震设防要求设备具有良好抗震性能。设计团队提出两个方案：方案A采用双母线双分段接线，户外AIS配电装置，主要设备包括户外支柱绝缘子、断路器、隔离开关等，设备布置需考虑盐雾防护措施（如增加爬电距离、涂覆防污闪涂料），但AIS设备整体抗震能力较弱，支柱绝缘子及断路器在强震下易发生断裂或位移；方案B采用3/2断路器接线，户外GIS配电装置，所有带电部件封闭在充有SF6气体的金属壳体内，完全隔离盐雾侵蚀，且GIS整体结构刚度高、重心低，抗震性能优异，但初期投资较高。设计团队需根据DL/T 5218-2012第5章（电气主接线）和第6章（配电装置），以及DL/T 5429-2009第6章（电网方案设计）的相关要求，对两个方案进行技术经济综合比选，并给出推荐方案。
  \item[A:] 现状诊断：本工程为500kV枢纽变电站，站址位于沿海高盐雾地区且地震设防烈度为8度，属于典型严酷环境条件。方案A采用双母线双分段接线与户外AIS配电装置，该配置在沿海高盐雾环境下存在绝缘子表面盐密积累导致的污闪风险，且AIS设备抗震性能较差，在8度地震作用下支柱绝缘子和断路器易发生断裂或位移，不符合DL/T 5218-2012第5章对枢纽变电站接线可靠性的要求，以及第6章对配电装置适应环境条件的要求。方案B采用3/2断路器接线与户外GIS配电装置，GIS将带电导体封闭在充有SF6气体的金属壳体内，完全隔离了盐雾侵蚀，且GIS整体结构刚度高、重心低，抗震性能显著优于AIS。然而，方案B初期投资较高，且GIS设备对密封工艺和安装环境要求严格，在沿海高盐雾条件下需特别关注气体泄漏和壳体腐蚀问题。总体而言，当前工程配置存在环境适应性不足与可靠性标准不匹配的核心矛盾。

多维冲突分析：本工程面临的核心冲突在于“高可靠性要求”与“严酷环境条件”之间的张力，以及“标准符合性”与“经济性”之间的权衡。具体而言，DL/T 5218-2012第5章要求500kV枢纽变电站应采用可靠性高的接线形式（如3/2断路器接线或双母线双分段），但方案A的双母线双分段接线在母线故障或检修时需切换母线，存在短时停电风险，且AIS设备在盐雾和地震作用下故障率显著升高，难以满足枢纽站“N-1”甚至“N-2”安全准则。DL/T 5429-2009第6章强调电网方案设计应满足DL 755《电力系统安全稳定导则》的要求，即电网结构应能承受单一故障而不损失负荷，而AIS在8度地震下的机械可靠性难以保证这一要求。另一方面，方案B的GIS虽能完美解决环境适应性问题，但投资较AIS高出约30\%-50\%，且全生命周期中需额外投入SF6气体监测与回收成本。冲突的本质在于：标准对安全性和可靠性的要求是刚性的，而工程投资受预算约束，必须在满足标准底线的前提下寻求最优经济方案。

多方案比选：从安全性角度，方案B的GIS完全密封，无污闪风险，且抗震性能经型式试验验证可承受8度地震，安全性显著优于方案A的AIS；方案A在盐雾和地震双重作用下，支柱绝缘子断裂、断路器拒动等事故概率大幅上升，存在系统性安全风险。从经济性角度，方案A初期投资低，但全生命周期成本因高故障率、频繁检修和停电损失而大幅增加；方案B初期投资高，但运维成本低、寿命长（GIS一般可达30年以上），全生命周期成本可控且更具经济性。从可行性角度，方案A的AIS设备技术成熟、安装简单，但需额外采取防污闪涂料、增加爬电距离、加强抗震加固等措施，实施复杂度不低；方案B的GIS需严格把控安装环境湿度与清洁度，且需配置SF6气体回收装置，对施工和运维要求较高，但在沿海地区已有大量成功案例，技术可行性充分。从标准符合性角度，方案B完全满足DL/T 5218-2012第5.2节对枢纽站接线可靠性的要求（3/2断路器接线可实现任一出线或断路器检修时不停电），以及第6.3节对配电装置适应环境的要求（GIS适用于高盐雾、高海拔、强震区）；方案A则需通过额外措施才能勉强达标，且存在标准符合性风险。综合比选，方案B在安全性、标准符合性和全生命周期经济性上全面优于方案A。

折中综合方案：推荐采用方案B（3/2断路器接线+户外GIS配电装置）作为最终实施方案。理由如下：首先，3/2断路器接线是500kV枢纽变电站的标准推荐接线形式，其可靠性满足DL/T 5218-2012第5章要求，且与DL/T 5429-2009第6章对电网安全稳定的要求高度一致。其次，GIS设备完全适应沿海高盐雾和8度地震设防的严酷环境，从根本上消除了污闪和抗震薄弱环节，避免了方案A需额外投入大量防污闪和抗震加固措施的成本与风险。第三，虽然初期投资较高，但全生命周期内因故障率极低、检修间隔长（GIS一般10-15年才需大修），综合经济性优于方案A。若因预算限制无法一次性实施全GIS方案，可考虑分阶段策略：近期先采用3/2断路器接线+混合配电装置（即关键线路采用GIS，次要线路采用AIS并加强防护），远期逐步替换为全GIS。但鉴于本工程为枢纽站，建议一次性投资到位，避免分期改造带来的停电损失和接口风险。

控制优先级链条：① 优先完成GIS设备的招标采购与技术协议签订，明确设备需满足8度抗震设防、盐雾等级III级（或以上）的型式试验要求，并约定SF6气体年泄漏率≤0.5\%的质保条款；② 同步开展GIS基础与构架设计，基础应进行抗震验算并设置隔震支座（如需要），确保与设备抗震性能匹配；③ 在土建施工阶段，优先完成GIS室（如采用户内GIS）或户外GIS基础的开挖与浇筑，严格控制基础沉降差；④ 设备安装前，对安装环境进行洁净度处理，确保湿度≤80\%、无粉尘，并配备SF6气体充注与回收装置；⑤ 安装过程中，严格执行GIS对接法兰的清洁与密封工艺，逐段进行真空检漏和微水测量；⑥ 投运前，完成全站3/2断路器接线的保护联调、断路器失灵保护及重合闸逻辑测试，确保接线可靠性；⑦ 投运后，建立GIS设备SF6气体压力与微水在线监测系统，并制定年度巡检与定期检漏计划，重点监测沿海盐雾环境下的壳体腐蚀情况。
  \item[KW:] DL/T 5429-2009、推荐方案、高地震烈度条件下、显著优于、断路器接线满足枢
  \item[标准:] DL/T 5429-2009, DL/T 5218-2012
  \item[条款溯源:] DL/T 5218-2012 第5章, DL/T 5218-2012 第6章, DL/T 5429-2009 第6章
  \item[主题:] 电气主接线
  \item[评分方式:] manual\_review
  \item[选项:] A \\ B
  \item[Rubric条款:] DL/T 5218-2012、DL/T 5429-2009、GB 38755-2019 第3.1.8条、DL/T 5429-2009 第6.2.3条
  \item[Rubric判据:] 方案A采用双母线双分段接线与户外AIS配电装置，该配置在沿海高盐雾环境下存在绝缘子表面盐密积累导致的污闪风险，且AIS设备抗震性能较差，在8度地震作用下支柱绝缘子和断路器易发生断裂或位移，不符合DL/T 5218-2012第5章对枢纽变电站接线可靠性的要求，以及第6章对配电装置适应环境条件的要求、然而，方案B初期投资较高，且GIS设备对密封工艺和安装环境要求严格，在沿海高盐雾条件下需特别关注气体泄漏和壳体腐蚀问题、总体而言，当前工程配置存在环境适应性不足与可靠性标准不匹配的核心矛盾、多维冲突分析：本工程面临的核心冲突在于“高可靠性要求”与“严酷环境条件”之间的张力，以及“标准符合性”与“经济性”之间的权衡、具体而言，DL/T 5218-2012第5章要求500kV枢纽变电站应采用可靠性高的接线形式（如3/2断路器接线或双母线双分段），但方案A的双母线双分段接线在母线故障或检修时需切换母线，存在短时停电风险，且AIS设备在盐雾和地震作用下故障率显著升高，难以满足枢纽站“N-1”甚至“N-2”安全准则
  \item[LJ判定:] 通过 (弱=0.40, 强=0.60, 差距=0.20)
\end{description}
\vspace{4pt}
\hrule
\vspace{6pt}

\noindent\textbf{L3-049} \hfill \textbf{综合评估} $|$ \textit{方案比选-变电站设计}
\begin{description}[leftmargin=1.5em,style=nextline,font=\normalfont]
  \item[Q:] 【场景】某220kV变电站位于沿海工业区，站址经实测盐雾密度等级为“重盐雾区”，且该区域规划为城市发展新区，周边将新建大型化工厂。设计阶段需确定电气主接线方案（第5章）及配电装置选型（第6章），同时需评估无功补偿配置（第7章）对系统安全稳定的影响。

【问题】根据DL/T 5218-2012第5章、第6章、第7章及GB 38755-2019第3.4节（无功平衡及补偿）、第4.2节（承受大扰动能力），请对比分析以下两个方案：方案A：采用双母线接线，屋外常规配电装置（AIS），在站内集中装设2组20Mvar并联电容器；方案B：采用双母线接线，屋内气体绝缘金属封闭开关设备（GIS），在站内集中装设2组20Mvar并联电容器+1组20Mvar并联电抗器。要求从环境适应性、系统无功平衡与电压稳定性、短路电流水平、设备抗污闪能力四个维度进行综合评判，并给出推荐方案。
  \item[A:] 现状诊断：该220kV变电站位于沿海重盐雾工业区，且周边规划新建大型化工厂，环境条件极为恶劣。当前方案A采用屋外常规配电装置（AIS），根据DL/T 5218-2012第3.0.11条要求，站址不宜设在严重盐雾地区，必要时应采取防污染措施，但AIS外绝缘爬电距离在重盐雾长期侵蚀下难以满足运行要求，存在严重污闪风险，不符合标准对恶劣环境下的设备选型规定。方案B采用屋内气体绝缘金属封闭开关设备（GIS），将高压带电体密封于金属壳内，从根本上隔离盐雾，满足标准要求。在无功补偿方面，方案A仅配置2组20Mvar并联电容器，缺乏吸收过剩无功的能力，而GB 38755-2019第3.4.1条要求无功补偿应满足分层分区平衡原则并具备动态调节能力，方案A在轻载或线路充电功率过大时易导致电压升高，存在电压稳定隐患。方案B增加1组20Mvar并联电抗器，可双向调节无功，更符合标准要求。

多维冲突分析：各方案与标准要求之间存在多重冲突。首先，环境适应性与设备选型冲突：方案A的AIS设备在重盐雾区无法满足DL/T 5218-2012第3.0.11条关于站址选择和防污染措施的要求，而方案B的GIS虽能解决此问题，但投资成本显著增加，与工程经济性要求产生矛盾。其次，无功补偿与电压稳定性冲突：GB 38755-2019第3.4.1条要求无功补偿应满足分层分区平衡并具备动态调节能力，方案A仅配置电容器，在轻载时无法吸收过剩无功，可能违反第4.2.1条关于承受大扰动后电压应在允许范围内的规定；方案B增加电抗器后虽能抑制工频过电压，但电抗器容量需与系统充电功率精确匹配，否则可能造成欠补偿或过补偿，引发新的电压问题。此外，短路电流控制冲突：DL/T 5218-2012第5章要求主接线设计应限制短路电流，方案A和B均采用双母线接线，短路电流水平相当，但方案B的GIS设备阻抗通常更小，可能使短路电流增大，需通过电抗器参数优化来平衡，这增加了设计复杂度。冲突的本质在于：在极端环境条件下，单一方案难以同时满足环境适应性、无功平衡、电压稳定和短路电流控制等多重标准要求，必须进行综合权衡。

多方案比选：从安全性维度看，方案B的GIS设备彻底隔离盐雾，消除了污闪风险，且具备双向无功调节能力，能更好地满足GB 38755-2019第4.2.1条关于承受大扰动后电压稳定的要求，安全性显著优于方案A。从经济性维度看，方案A的AIS设备初始投资较低，但长期运行中因污闪导致的维护和停电损失风险高；方案B的GIS设备投资高，但运行可靠性高、维护成本低，全生命周期经济性可能更优。从可行性维度看，方案A在重盐雾区需额外采取防污闪措施（如加强绝缘子清扫、涂覆防污涂料等），运行管理复杂且效果有限；方案B的GIS技术成熟，安装施工方便，但需考虑屋内配电装置的土建施工和通风散热问题。从标准符合性维度看，方案A不符合DL/T 5218-2012第3.0.11条关于站址选择的要求，且无功补偿能力不足，违反GB 38755-2019第3.4.1条；方案B完全满足上述标准要求，同时通过电抗器配置符合第3.4.2条“无功补偿应具有足够的调节能力”的规定。综合对比，方案B在安全性、标准符合性和长期可行性上全面优于方案A，经济性虽初始投资高，但全生命周期成本可控。

折中综合方案：推荐采用方案B，即采用双母线接线、屋内GIS配电装置，并配置2组20Mvar并联电容器和1组20Mvar并联电抗器。该方案通过GIS彻底解决重盐雾区的污闪问题，满足DL/T 5218-2012第3.0.11条要求；通过并联电抗器实现无功双向调节，在轻载时吸收过剩无功、抑制工频过电压，在重载时提供无功支撑，符合GB 38755-2019第3.4.1条和第4.2.1条关于无功平衡及承受大扰动能力的规定。若因投资限制无法一步到位，可考虑分阶段实施：初期先建设GIS和电容器组，确保环境适应性和基本无功补偿；后期根据系统实际运行情况，增建电抗器组。在工程实施中，需进一步校核电抗器容量与系统充电功率的匹配性，避免欠补偿或过补偿；同时按照DL/T 5218-2012第7.2节要求配置自动投切装置，实现无功补偿的自动调节，确保电压稳定。

控制优先级链条：① 优先完成GIS设备的选型与采购，确保设备满足重盐雾区的密封和防腐要求，并按照DL/T 5218-2012第3.0.11条进行站址环境复核；② 同步开展无功补偿容量的详细计算，根据系统充电功率和负荷特性，精确确定电抗器和电容器的容量配比，确保符合GB 38755-2019第3.4.1条分层分区平衡原则；③ 设计并安装自动投切装置，按照DL/T 5218-2012第7.2节要求，实现电容器和电抗器的自动调节，满足动态无功补偿需求；④ 进行短路电流校核计算，根据GIS设备阻抗和电抗器参数，验证短路电流是否在断路器额定开断能力范围内，必要时调整主接线或增加限流措施；⑤ 制定运行维护规程，明确GIS设备的密封检测周期、电抗器和电容器的投切策略，以及电压越限时的应急处置流程，确保长期安全稳定运行。
  \item[KW:] DL/T 5218-2012、GB 38755-2019、分层分区、无功补偿、无功平衡
  \item[标准:] GB 38755-2019, DL/T 5218-2012
  \item[条款溯源:] GB 38755-2019 第3.4.1条, DL/T 5218-2012 第3.0.11条
  \item[主题:] 电压稳定
  \item[评分方式:] manual\_review
  \item[Rubric条款:] DL/T 5218-2012、GB 38755-2019、DL/T 5218-2012 第3.0.11条、GB 38755-2019 第3.4.1条、GB 38755-2019 第4.2.1条、DL/T 5218-2012 第5章、GB 38755-2019 第3.4.2条、DL/T 5218-2012 第7.2节
  \item[Rubric判据:] 11条要求，站址不宜设在严重盐雾地区，必要时应采取防污染措施，但AIS外绝缘爬电距离在重盐雾长期侵蚀下难以满足运行要求，存在严重污闪风险，不符合标准对恶劣环境下的设备选型规定、方案B采用屋内气体绝缘金属封闭开关设备（GIS），将高压带电体密封于金属壳内，从根本上隔离盐雾，满足标准要求、1条要求无功补偿应满足分层分区平衡原则并具备动态调节能力，方案A在轻载或线路充电功率过大时易导致电压升高，存在电压稳定隐患、方案B增加1组20Mvar并联电抗器，可双向调节无功，更符合标准要求、多维冲突分析：各方案与标准要求之间存在多重冲突
  \item[LJ判定:] 通过 (弱=0.80, 强=1.00, 差距=0.20)
\end{description}
\vspace{4pt}
\hrule
\vspace{6pt}

\noindent\textbf{L3-009} \hfill \textbf{综合评估} $|$ \textit{综合决策-短路电流控制}
\begin{description}[leftmargin=1.5em,style=nextline,font=\normalfont]
  \item[Q:] 【场景】某区域电网现有500kV/220kV电磁环网运行，500kV主网架由3座500kV变电站（A、B、C）构成，每站主变容量2×750MVA，220kV电网由8座220kV变电站组成，线路总长约200公里。目前短路电流水平高达52kA，超过断路器遮断容量上限50kA，计划解环运行以控制短路电流。解环后，220kV电网分为两个片区，每个片区通过一台500kV联络变压器与主网相连。两个片区之间无220kV联络线。设计水平年2028年。约束条件：解环后每个片区短路电流需控制在45kA以下；220kV线路走廊紧张，新建线路需穿越城区，环保审批严格；联络变压器容量限制为600MVA。决策目标：在满足短路电流限制、走廊和环保约束的前提下，综合评估两种解环方案对供电可靠性、运行灵活性、投资成本及事故备用能力的影响，选出最优方案。

【问题】请对以下两种解环方案进行综合评估。
方案A：在500kV变电站A处解环，将220kV母线分裂为两段，分别供两个片区，每段母线通过一台联络变与500kV主网相连。
方案B：在500kV变电站B处解环，同时新建1回220kV联络线连接两个片区，作为事故备用。
  \item[A:] 现状诊断：现有电磁环网导致500kV/220kV短路电流均超标（超过50kA），需解环运行。但解环后，两个220kV片区将失去相互支援能力，供电可靠性下降。需同时满足GB 38755-2019对“合理的电网结构”（第3.2节）和DL/T 5429-2009对“短路电流控制”（第6.1.8条）及“分层分区”（第6.1.5条）的要求。

多维冲突分析：
1. 短路电流与供电可靠性冲突：方案A解环彻底，短路电流可降至40kA以下，但两个片区完全独立，任一联络变故障将导致该片区全停（违反GB 38755-2019第2.3条N-1原则）。方案B保留联络线，可靠性高但短路电流降幅有限。
2. 经济性与灵活性冲突：方案A投资省（无需新建线路），但运行灵活性差。方案B需新建220kV线路（投资约5000万元），但提供事故支援能力。
3. 电压控制与无功平衡冲突：解环后，各片区无功平衡难度加大（GB 38755-2019第3.4条），方案A需在各片区新增无功补偿装置。

多方案比选：
- 方案A：优点：①短路电流可控制在40kA以内（符合DL/T 5429-2009第6.1.8条）；②投资省。缺点：①不满足GB 38755-2019第3.2节“消除薄弱环节”要求，联络变N-1故障将导致片区全停；②电压控制困难（需新增4组60Mvar电容器）。
- 方案B：优点：①满足N-1原则（联络变故障后可通过联络线转供负荷）；②短路电流可控制在45kA以内（仍可接受）；③电压控制相对容易。缺点：①投资较高；②需增加保护及安全自动装置。

折中综合方案：推荐方案B，但优化为：①新建的220kV联络线配置自动解列装置（当短路电流超过阈值时自动断开，兼顾安全与可靠性）；②在每个片区配置1台调相机（±100Mvar），解决解环后的电压稳定问题（GB 38755-2019第5.6节）；③联络变压器采用高阻抗变压器，进一步限制短路电流。

控制优先级链条：
第一优先级（短路电流）：不超过设备耐受值（DL/T 5429-2009第6.1.8条）→ 第二优先级（N-1原则）：任一元件故障不损失负荷（GB 38755-2019第2.3条）→ 第三优先级（电网结构）：符合分层分区原则（DL/T 5429-2009第6.1.5条）→ 第四优先级（电压稳定）：配置无功补偿（GB 38755-2019第3.4条）→ 第五优先级（经济性）：在满足前三项前提下优化投资。
  \item[KW:] GB 38755-2019 第2.3条、DL/T 5429-2009 第6.1.8条、电磁环网解环、短路电流控制、N-1原则
  \item[标准:] GB 38755-2019, DL/T 5429-2009
  \item[条款溯源:] GB 38755-2019 第2.3条, DL/T 5429-2009 第6.1.8条, DL/T 5429-2009 第6.1.5条, GB 38755-2019 第3.2节, GB 38755-2019 第3.4条
  \item[主题:] 短路电流
  \item[评分方式:] manual\_review
  \item[Rubric条款:] GB 38755-2019、DL/T 5429-2009、GB 38755-2019 第3.2节、DL/T 5429-2009 第6.1.8条、DL/T 5429-2009 第6.1.5条、GB 38755-2019 第2.3条、GB 38755-2019 第3.4条、GB 38755-2019 第5.6节
  \item[Rubric判据:] 现状诊断：现有电磁环网导致500kV/220kV短路电流均超标（超过50kA），需解环运行、需同时满足GB 38755-2019对“合理的电网结构”（第3、经济性与灵活性冲突：方案A投资省（无需新建线路），但运行灵活性差、方案B需新建220kV线路（投资约5000万元），但提供事故支援能力、4条），方案A需在各片区新增无功补偿装置
  \item[LJ判定:] 需重写 (弱=0.20, 强=0.20, 差距=0.00)
  \item[LJ反馈:] L3区分度不足：差距=0.00。请增强方案间的技术冲突或增加工程决策的难度。
\end{description}
\vspace{4pt}
\hrule
\vspace{6pt}

\noindent\textbf{L3-001} \hfill \textbf{综合评估} $|$ \textit{方案比选-调峰与备用}
\begin{description}[leftmargin=1.5em,style=nextline,font=\normalfont]
  \item[Q:] 【场景】某区域电网规划水平年（2028年）最大负荷为25000MW，现有电源装机容量为28000MW，其中火电占70\%，水电占20\%，新能源（风电、光伏）占10\%。系统负荷备用率按3\%考虑，事故备用率按10\%考虑。规划需新增电源以满足负荷增长和备用需求，同时需满足电网安全稳定要求。现有电网结构为500kV环网，但存在部分220kV长距离输电线路，且新能源接入点附近无功补偿不足。

【问题】请基于GB 38755-2019和DL/T 5429-2009，对以下两种电源建设方案进行综合评估，并给出推荐方案。
方案A：新建2×600MW超超临界燃煤火电机组，接入500kV枢纽变电站，同步建设300Mvar调相机。
方案B：新建1000MW集中式光伏电站+500MW/1000MWh电化学储能电站，通过220kV线路接入系统，配套建设200Mvar SVG动态无功补偿装置。
  \item[A:] 现状诊断：当前系统总装机28000MW，最大负荷25000MW，备用容量仅3000MW（10.7\%），低于DL/T 5429-2009第5.2.3条要求的15\%-20\%（即3750MW-5000MW），且事故备用容量（10\%）虽满足下限，但未明确是否大于系统一台最大单机容量。同时，新能源占比10\%，但无功补偿不足，存在电压稳定风险，不符合GB 38755-2019第3.4条关于无功平衡及补偿的要求。

多维冲突分析：
1. 安全稳定维度：方案A新增火电机组（常规电源）可提供惯量支撑和电压支撑，有利于维持系统暂态功角稳定和电压稳定（GB 38755-2019第4.2条）；方案B新能源+储能虽可提供快速响应，但光伏无转动惯量，大规模接入可能降低系统惯性，增加频率稳定风险（GB 38755-2019第5.7条）。
2. 经济性维度：方案A初投资高，燃料成本受煤价波动影响大；方案B光伏成本持续下降，储能可参与调峰调频，但需考虑储能寿命和更换成本。
3. 环保与政策维度：方案A碳排放高，不符合双碳目标；方案B清洁低碳，符合国家能源转型方向（DL/T 5429-2009第5.1.3条第5款）。
4. 电网适应性维度：方案A接入500kV枢纽站，对网架结构影响小；方案B通过220kV线路接入，可能加重已有220kV长距离线路的输送压力，需校核N-1原则（GB 38755-2019第2.3条）。

多方案比选：
- 方案A优势：提供系统惯量和短路容量，增强系统强度；调相机可提供动态无功支撑，改善电压稳定性（GB 38755-2019第3.4条）；劣势：建设周期长（4-5年），碳排放高，燃料供应存在不确定性。
- 方案B优势：建设周期短（1-2年），清洁低碳，储能可提供快速功率支撑，参与调峰（DL/T 5429-2009第5.3.1条）；劣势：无惯量，需依赖系统其他电源提供惯量支撑；储能容量有限，持续支撑能力弱；220kV接入需加强网架。

折中综合方案：推荐采用方案B为主，但需补充以下措施：
1. 在新能源接入点加装同步调相机或构网型储能，提升系统惯量和短路比（GB 38755-2019第3.3条、第3.4条）。
2. 同步建设1回500kV线路将220kV区域与主网加强联络，满足N-1原则（GB 38755-2019第2.3条）。
3. 保留部分火电机组作为旋转备用，确保系统频率稳定（GB 38755-2019第5.7条）。

控制优先级链条：
1. 首先满足DL/T 5429-2009第5.2.3条备用容量要求（总备用≥15\%），通过方案B+保留火电备用实现。
2. 其次满足GB 38755-2019第4.2条大扰动安全稳定标准（暂态功角稳定、电压稳定），通过加装调相机/构网型储能实现。
3. 然后满足GB 38755-2019第3.4条无功分层分区平衡要求，通过SVG+调相机实现。
4. 最后满足DL/T 5429-2009第6.1.4条电网结构安全要求，通过加强500kV联络线实现。
  \item[KW:] 备用容量、无功补偿、惯量支撑、N-1原则、暂态稳定
  \item[标准:] GB 38755-2019, DL/T 5429-2009
  \item[条款溯源:] GB 38755-2019 第3.4条, GB 38755-2019 第4.2条, GB 38755-2019 第5.7条, GB 38755-2019 第2.3条, DL/T 5429-2009 第5.2.3条, DL/T 5429-2009 第5.1.3条第5款
  \item[主题:] 电压稳定
  \item[评分方式:] manual\_review
  \item[Rubric条款:] DL/T 5429-2009、GB 38755-2019、DL/T 5429-2009 第5.2.3条、GB 38755-2019 第3.4条、GB 38755-2019 第4.2条、GB 38755-2019 第5.7条、DL/T 5429-2009 第5.1.3条第5款、GB 38755-2019 第2.3条
  \item[Rubric判据:] 7\%），低于DL/T 5429-2009第5、3条要求的15\%-20\%（即3750MW-5000MW），且事故备用容量（10\%）虽满足下限，但未明确是否大于系统一台最大单机容量、同时，新能源占比10\%，但无功补偿不足，存在电压稳定风险，不符合GB 38755-2019第3、方案B新能源+储能虽可提供快速响应，但光伏无转动惯量，大规模接入可能降低系统惯性，增加频率稳定风险（GB 38755-2019第5、方案B光伏成本持续下降，储能可参与调峰调频，但需考虑储能寿命和更换成本
  \item[LJ判定:] 通过 (弱=0.40, 强=0.60, 差距=0.20)
\end{description}
\vspace{4pt}
\hrule
\vspace{6pt}

\noindent\textbf{L3-041} \hfill \textbf{综合评估} $|$ \textit{多标准协调-无功配置}
\begin{description}[leftmargin=1.5em,style=nextline,font=\normalfont]
  \item[Q:] 【场景】某750kV变电站无功补偿设计，该变电站位于负荷中心附近，主变容量为2×1500MVA，额定电压750/330kV，330kV侧出线6回，每回线路长度约80km，输送容量约800MW，短路电流水平为40kA。根据电网规划，该区域新能源装机容量较大，无功电压波动频繁。现提出两种无功补偿方案：方案A采用集中式并联电容器组，在750kV母线侧集中安装4组60Mvar电容器，总容量240Mvar；方案B采用分散式并联电抗器+电容器组，在330kV母线侧分散安装8组30Mvar电容器和4组30Mvar电抗器，总容量360Mvar（电容器240Mvar，电抗器120Mvar）。需依据DL/T 5218-2012第7章无功补偿和GB 38755-2019第3.4节无功平衡及补偿进行比选。

【问题】根据DL/T 5218-2012第7章和GB 38755-2019第3.4节，分析两种无功补偿方案在电压调节能力、无功分层分区平衡及经济性方面的差异，并给出推荐方案。
  \item[A:] 现状诊断：当前750kV变电站无功补偿设计面临两种方案的选择。方案A采用集中式并联电容器组，根据DL/T 5218-2012第7章，集中补偿便于集中管理和维护，投资较低，但该方案可能导致无功功率长距离输送，增加网损，且电压调节灵活性较差，难以适应局部无功需求变化。方案B采用分散式并联电抗器+电容器组，DL/T 5218-2012第7章支持分散补偿，可有效抑制工频过电压和操作过电压，但投资较高、占地较大。依据GB 38755-2019第3.4节关于无功补偿应分层分区、就地平衡的要求，方案A的集中补偿模式存在违反该标准原则的风险，因为其无法实现无功的精确分层分区平衡，可能造成无功在电网中长距离流动，不利于电压稳定。

多维冲突分析：本场景的核心冲突在于经济性与标准符合性及技术性能之间的矛盾。方案A（集中式）在投资经济性上具有明显优势，符合DL/T 5218-2012中关于集中管理、降低初始投资的要求，但其无功补偿方式与GB 38755-2019第3.4节强调的“分层分区、就地平衡”原则相冲突。集中补偿无法针对局部无功需求进行精准调节，导致无功功率在长距离传输中产生额外损耗，并可能引发电压波动，违反了安全稳定导则对电压质量的要求。方案B（分散式）虽然完全符合GB 38755-2019第3.4节的要求，能实现无功就地平衡和精细化的电压调节，但其高昂的投资成本和较大的占地面积又与工程经济性要求产生冲突。冲突的本质在于：在超高压（750kV）变电站这一特定场景下，系统对电压稳定性和无功平衡的刚性技术要求（由GB 38755-2019规定）与工程投资的经济性约束（由DL/T 5218-2012部分体现）之间存在根本性矛盾。

多方案比选：从安全性角度，方案B显著优于方案A。方案B通过分散配置电抗器和电容器，能有效吸收线路充电功率、抑制工频过电压，并实现无功就地平衡，大幅提升电压稳定性，完全符合GB 38755-2019第3.4节的安全要求；而方案A因集中补偿导致无功长距离输送，电压调节灵活性差，在750kV这种高电压等级下存在电压失稳风险。从经济性角度，方案A优于方案B，其初始投资低、占地小、维护集中；方案B投资高、占地大。从可行性角度，两种方案在技术上都可行，但方案B更符合750kV变电站对无功补偿的精细化要求。从标准符合性角度，方案B完全满足GB 38755-2019第3.4节关于无功分层分区平衡的要求，也符合DL/T 5218-2012第7章对超高压变电站推荐分散补偿的导向；方案A则与GB 38755-2019第3.4节的原则相悖。综合来看，方案B在安全性和标准符合性上具有压倒性优势，尽管经济性较差，但对于750kV这一关键枢纽变电站，安全稳定运行是首要目标。

折中综合方案：推荐采用方案B（分散式并联电抗器+电容器组）作为最终实施方案。主要理由如下：第一，750kV变电站作为电网的骨干节点，其线路充电功率巨大，电压控制要求极高，方案B的分散式配置能够精确实现无功分层分区平衡，有效抑制过电压，这是保障电网安全稳定运行的根本要求，完全符合GB 38755-2019第3.4节的强制性规定。第二，DL/T 5218-2012第7章也明确推荐在超高压变电站采用分散补偿方式以提高电压调节能力，方案B与该标准的技术导向一致。虽然方案B投资较高，但考虑到750kV变电站的安全重要性，该投资是必要且合理的。为缓解经济性压力，可考虑在初步设计阶段优化设备选型和布局，例如采用紧凑型组合式设备以减少占地，或通过招标采购控制设备成本，但不应牺牲分散补偿的核心原则。

控制优先级链条：① 优先进行详细的系统无功电压计算分析，明确750kV变电站及其周边电网在不同运行方式下的无功需求、电压波动范围及线路充电功率大小，为分散式补偿装置的容量和布点提供精确依据。② 根据计算结果，确定并联电抗器和电容器组的分散配置方案，包括各补偿点的具体容量、分组投切策略及安装位置，确保满足GB 38755-2019第3.4节分层分区平衡的要求。③ 开展设备选型与招标工作，优先选择技术成熟、可靠性高、占地小的设备（如干式空心电抗器、集合式电容器组），在满足技术性能的前提下通过竞争性采购控制投资成本。④ 制定详细的分阶段实施计划，优先安装对电压稳定影响最大的关键节点的补偿装置，例如先安装主变压器低压侧的电抗器以吸收充电功率，再逐步配置各出线侧的电容器组。⑤ 在工程实施过程中，同步建设配套的自动电压控制（AVC）系统，确保分散式补偿装置能够根据电网运行状态自动、协调地进行投切，实现无功电压的闭环控制。
  \item[KW:] DL/T 5218-2012、GB 38755-2019、分层分区、就地平衡、无功补偿
  \item[标准:] GB 38755-2019, DL/T 5218-2012
  \item[条款溯源:] GB 38755-2019 第3.2.1条
  \item[主题:] 无功平衡
  \item[评分方式:] manual\_review
  \item[Rubric条款:] DL/T 5218-2012、GB 38755-2019、GB 38755-2019 第3.4节、DL/T 5218-2012 第7章
  \item[Rubric判据:] 方案A采用集中式并联电容器组，根据DL/T 5218-2012第7章，集中补偿便于集中管理和维护，投资较低，但该方案可能导致无功功率长距离输送，增加网损，且电压调节灵活性较差，难以适应局部无功需求变化、4节关于无功补偿应分层分区、就地平衡的要求，方案A的集中补偿模式存在违反该标准原则的风险，因为其无法实现无功的精确分层分区平衡，可能造成无功在电网中长距离流动，不利于电压稳定、多维冲突分析：本场景的核心冲突在于经济性与标准符合性及技术性能之间的矛盾、方案A（集中式）在投资经济性上具有明显优势，符合DL/T 5218-2012中关于集中管理、降低初始投资的要求，但其无功补偿方式与GB 38755-2019第3、集中补偿无法针对局部无功需求进行精准调节，导致无功功率在长距离传输中产生额外损耗，并可能引发电压波动，违反了安全稳定导则对电压质量的要求
  \item[LJ判定:] 通过 (弱=0.80, 强=1.00, 差距=0.20)
\end{description}
\vspace{4pt}
\hrule
\vspace{6pt}

\noindent\textbf{L3-094} \hfill \textbf{综合评估} $|$ \textit{综合决策-继电保护}
\begin{description}[leftmargin=1.5em,style=nextline,font=\normalfont]
  \item[Q:] 【场景】某地区新建一座110kV变电站，主变容量为2×50MVA，110kV侧采用单母线分段接线，出线4回，每回线路长度约15km，最大输送容量40MW，系统短路电流水平为12kA。该变电站需设计继电保护及监控系统，现提出两个方案：方案A采用常规电磁式继电保护+就地监控屏，保护装置采用电磁型继电器，动作速度约40-60ms，无自检功能，故障率较高，就地监控屏依赖人工巡检，信息上传延迟大，难以满足现代电网对快速故障切除和广域监控的需求；方案B采用微机综合自动化系统，保护、测控、通信一体化，保护动作速度可达20-30ms，具备自诊断和远方通信功能，但依赖通信网络，存在网络攻击风险。需根据DL/T 5218-2012第8章（电气二次/继电保护/监控）和DL/T 5429-2009第6章（系统及电气二次）进行评价。

【问题】依据DL/T 5218-2012第6.4节（计算机监控与二次接线）和DL/T 5429-2009第6.1.4条（安全稳定导则），比较方案A与方案B在保护动作速度、系统可靠性及对电网安全稳定支撑方面的差异，并推荐方案。
  \item[A:] 现状诊断：当前110kV变电站设计面临方案A（常规电磁式继电保护+就地监控屏）与方案B（微机综合自动化系统）的选型问题。方案A采用电磁式继电器，保护动作速度约40-60ms，无自检功能，故障率较高，就地监控屏依赖人工巡检，信息上传延迟大，难以满足现代电网对快速故障切除和广域监控的需求。方案B采用微机综合自动化，保护动作速度可达20-30ms，具备自诊断和远方通信功能，但依赖通信网络，存在网络攻击风险。依据DL/T 5218-2012第6.4.1条，计算机监控系统应满足安全可靠、经济适用原则，方案A在安全可靠性上存在明显不足；依据DL/T 5429-2009第6.1.4条，系统设计需满足DL 755安全稳定导则，方案A的慢速保护可能影响系统暂态稳定，不符合标准要求。

多维冲突分析：各方案与标准要求之间存在多重冲突。首先，方案A的电磁式保护动作速度（40-60ms）与DL 755对暂态稳定要求的快速故障切除需求相冲突，慢速保护可能导致系统失稳；而方案B虽满足速度要求，但其对通信网络的依赖与DL/T 5218-2012第6.4.3条要求的远方控制功能存在潜在冲突——通信网络故障可能导致远方控制失效。其次，方案A的低投资成本与DL/T 5218-2012第6.4.1条“安全可靠”原则冲突，低成本以牺牲可靠性为代价；方案B的高投资（约高40\%）与“经济适用”原则存在表面冲突，但全生命周期效益可能更优。冲突的本质在于：传统技术（方案A）无法同时满足现代电网对快速性、可靠性和智能化监控的多重要求，而先进技术（方案B）虽能达标，却引入了新的风险（如网络安全）和成本压力。

多方案比选：从安全性看，方案B的快速保护（20-30ms）可有效维持系统暂态稳定，优于方案A的40-60ms；方案B的自诊断功能可减少拒动概率，但通信网络故障可能导致功能丧失，而方案A无此风险但整体可靠性更低。从经济性看，方案A初始投资较低，但方案B全生命周期效益显著，包括减少停电损失、降低运维成本（远方监控替代人工巡检）以及便于未来智能化升级。从可行性看，方案B符合DL/T 5218-2012第6.4.1条的技术发展方向，且便于集成安全自动装置（如低频减载），而方案A难以实现此类功能。从标准符合性看，方案B完全满足DL/T 5429-2009第6.1.4条对安全稳定导则的要求，方案A则存在明显差距。综合而言，方案B在安全性、标准符合性和长期可行性上全面占优，经济性虽初始较高但全生命周期更优。

折中综合方案：推荐采用方案B（微机综合自动化系统），理由如下：DL/T 5218-2012第6.4.1条要求监控系统安全可靠，方案B符合技术发展方向；DL/T 5429-2009第6.1.4条要求满足安全稳定导则，方案B的快速保护可有效维持系统稳定。虽然投资较高，但全生命周期效益显著，且便于未来智能化升级。为平衡经济性与风险，可采取分阶段实施策略：第一阶段优先部署微机保护及测控单元，实现快速故障切除和本地监控；第二阶段逐步完善通信网络和远方控制功能，并同步加强网络安全防护（如加密、防火墙、入侵检测），以降低网络攻击风险。此方案既满足标准要求，又通过分阶段投入缓解一次性投资压力。

控制优先级链条：①第一阶段：优先采购并安装微机保护装置和测控单元，替换电磁式继电器，实现保护动作速度从40-60ms提升至20-30ms，满足DL 755暂态稳定要求；②同步部署本地监控系统，实现故障信息实时记录和自诊断功能，减少拒动概率；③第二阶段：建设通信网络（光纤或以太网），连接各保护测控单元至站控层，实现远方通信和远方控制功能，符合DL/T 5218-2012第6.4.3条要求；④同步实施网络安全防护措施，包括网络隔离、身份认证、数据加密和入侵检测系统，防范网络攻击风险；⑤第三阶段：集成安全自动装置（如低频减载、备用电源自投），提升电网安全稳定支撑能力；⑥最终完成全站智能化升级，接入调度自动化系统，实现广域监控和协同控制。
  \item[KW:] DL/T 5218-2012、DL/T 5429-2009、低频减载、安全自动装置、分层分区
  \item[标准:] DL/T 5429-2009, DL/T 5218-2012
  \item[条款溯源:] DL/T 5218-2012 第6.4.1条, DL/T 5429-2009 第6.1.4条
  \item[主题:] 暂态功角稳定
  \item[评分方式:] manual\_review
  \item[Rubric条款:] DL/T 5218-2012、DL/T 5429-2009、DL/T 5218-2012 第6.4.1条、DL/T 5429-2009 第6.1.4条、DL 755、DL/T 5218-2012 第6.4.3条
  \item[Rubric判据:] 方案A采用电磁式继电器，保护动作速度约40-60ms，无自检功能，故障率较高，就地监控屏依赖人工巡检，信息上传延迟大，难以满足现代电网对快速故障切除和广域监控的需求、1条，计算机监控系统应满足安全可靠、经济适用原则，方案A在安全可靠性上存在明显不足、4条，系统设计需满足DL 755安全稳定导则，方案A的慢速保护可能影响系统暂态稳定，不符合标准要求、多维冲突分析：各方案与标准要求之间存在多重冲突、首先，方案A的电磁式保护动作速度（40-60ms）与DL 755对暂态稳定要求的快速故障切除需求相冲突，慢速保护可能导致系统失稳
  \item[LJ判定:] 通过 (弱=0.40, 强=0.60, 差距=0.20)
\end{description}
\vspace{4pt}
\hrule
\vspace{6pt}

\noindent\textbf{L3-058} \hfill \textbf{综合评估} $|$ \textit{方案比选-变电站设计}
\begin{description}[leftmargin=1.5em,style=nextline,font=\normalfont]
  \item[Q:] 【场景】某330kV变电站，其330kV配电装置采用一个半断路器接线，初期仅形成两个完整串。该变电站位于区域电网的负荷中心，通过两回330kV线路与主网相连，线路长度分别为80km和120km，单回线路最大输送容量为800MW。站内安装两台容量为360MVA的主变压器，分别接入两个完整串中。系统短路电流水平为40kA。系统要求该站具备承受大扰动（如一台主变跳闸）后保持暂态功角稳定的能力。在初期接线中，隔离开关的配置存在两种方案：方案A为在各元件出口处（包括线路、变压器及母线侧）均装设隔离开关；方案B为仅在线路和变压器出口装设隔离开关，母线侧不装。
  \item[A:] 现状诊断：该330kV变电站采用一个半断路器接线，初期仅形成两个完整串，系统要求具备承受大扰动（如一台主变跳闸）后保持暂态功角稳定的能力。当前工程配置面临隔离开关选型问题：方案A在各元件出口处均装设隔离开关，方案B仅在线路和变压器出口装设，母线侧不装。根据DL/T 5218-2012第5.1.8条，初期接线“各元件出口处宜装设隔离开关”，方案B在母线侧未装设，不符合该条款的推荐性要求；同时，GB 38755-2019第4.2条要求电力系统承受大扰动时应保持稳定，第5.4条要求进行暂态功角稳定计算分析，方案B因母线侧无隔离开关，检修困难可能导致故障扩大或恢复时间延长，不利于满足暂态稳定要求，存在安全设计隐患。

多维冲突分析：各方案与标准要求之间存在多重冲突。首先，DL/T 5218-2012第5.1.8条建议“各元件出口处宜装设隔离开关”，但方案B在母线侧不装，直接违背了该条款的推荐性要求，而方案A则完全符合。其次，运行灵活性与经济性之间存在冲突：方案A虽灵活性高、便于检修隔离，但增加了设备投资和占地面积；方案B虽节省设备，但母线侧无隔离开关导致母线或断路器检修时需停运整个串，灵活性差。再者，继电保护配合方面，方案A利用隔离开关辅助触点提供明确隔离状态，便于保护逻辑判断；方案B依赖断路器状态，保护配合复杂，增加了误动风险。最后，暂态功角稳定方面，GB 38755-2019第5.4条要求计算分析大扰动后稳定性，方案A可快速隔离故障元件，减少故障持续时间，有利于暂态稳定；方案B因检修困难可能导致故障扩大或恢复时间延长，对暂态稳定不利。冲突的本质在于：在满足标准安全性与运行灵活性要求的前提下，如何平衡设备投资与系统可靠性之间的矛盾。

多方案比选：从安全性角度，方案A因隔离开关全装，可快速隔离故障元件，减少故障持续时间，有利于暂态功角稳定，符合GB 38755-2019第4.2条和第5.4条要求；方案B因母线侧无隔离开关，检修困难可能导致故障扩大，安全性较低。从经济性角度，方案A增加了设备投资和占地面积，方案B节省了设备成本，但需考虑因灵活性差导致的停电损失和运维成本。从可行性角度，方案A技术成熟，施工和运维经验丰富，易于实施；方案B虽减少设备，但保护配合复杂，对运维人员要求高，且母线侧无隔离开关可能违反DL/T 5218-2012第5.1.8条建议。从标准符合性角度，方案A完全符合DL/T 5218-2012第5.1.8条“各元件出口处宜装设隔离开关”的要求，方案B则不符合。综合来看，方案A在安全性、标准符合性和可行性上均优于方案B，经济性上的劣势可通过长期运行可靠性收益弥补。

折中综合方案：推荐采用方案A，即在各元件出口处均装设隔离开关。理由如下：首先，方案A完全符合DL/T 5218-2012第5.1.8条对初期接线“各元件出口处宜装设隔离开关”的建议，避免了标准合规风险。其次，在暂态功角稳定方面，方案A可快速隔离故障元件，减少故障持续时间，满足GB 38755-2019第4.2条和第5.4条对承受大扰动后保持稳定的要求。第三，方案A运行灵活性高，便于检修和隔离，有利于提高供电可靠性和运维效率。虽然初期投资较高，但考虑到该站为330kV枢纽变电站，且系统要求具备承受大扰动后的暂态稳定能力，安全性应优先于经济性。若因占地或投资限制需优化，可考虑在初期仅对关键元件（如主变和线路出口）装设隔离开关，但母线侧必须保留隔离开关，以确保检修灵活性，这实质上是方案A的简化实施。

控制优先级链条：① 优先实施方案A，在各元件出口处（包括线路、变压器和母线侧）均装设隔离开关，确保完全符合DL/T 5218-2012第5.1.8条要求。② 在隔离开关选型时，优先采用具有辅助触点的型号，以便为继电保护提供明确的隔离状态信号，简化保护逻辑判断。③ 在施工阶段，同步完成隔离开关与断路器、保护装置的联调测试，确保辅助触点信号准确可靠。④ 在投运前，根据GB 38755-2019第5.4条要求，进行暂态功角稳定计算分析，验证方案A在承受主变跳闸等大扰动后的稳定性。⑤ 在运维阶段，制定基于隔离开关的检修隔离操作流程，充分利用其灵活性，缩短故障隔离和恢复时间，进一步提升暂态稳定裕度。
  \item[KW:] DL/T 5218-2012、GB 38755-2019、暂态功角稳定、隔离开关
  \item[标准:] GB 38755-2019, DL/T 5218-2012
  \item[条款溯源:] GB 38755-2019 第5.4条, DL/T 5218-2012 第5.1.8条
  \item[主题:] 电气主接线
  \item[评分方式:] manual\_review
  \item[Rubric条款:] DL/T 5218-2012、GB 38755-2019、DL/T 5218-2012 第5.1.8条、GB 38755-2019 第4.2条、GB 38755-2019 第5.4条
  \item[Rubric判据:] 现状诊断：该330kV变电站采用一个半断路器接线，初期仅形成两个完整串，系统要求具备承受大扰动（如一台主变跳闸）后保持暂态功角稳定的能力、8条，初期接线“各元件出口处宜装设隔离开关”，方案B在母线侧未装设，不符合该条款的推荐性要求、2条要求电力系统承受大扰动时应保持稳定，第5、4条要求进行暂态功角稳定计算分析，方案B因母线侧无隔离开关，检修困难可能导致故障扩大或恢复时间延长，不利于满足暂态稳定要求，存在安全设计隐患、多维冲突分析：各方案与标准要求之间存在多重冲突
  \item[LJ判定:] 通过 (弱=1.00, 强=1.00, 差距=0.00)
\end{description}
\vspace{4pt}
\hrule
\vspace{6pt}

\noindent\textbf{L3-014} \hfill \textbf{综合评估} $|$ \textit{多标准协调-无功配置}
\begin{description}[leftmargin=1.5em,style=nextline,font=\normalfont]
  \item[Q:] 【场景】某区域电网规划建设一座大型火电厂（4×600MW），通过500kV线路接入主网。该区域同时有大规模风电基地（2000MW）接入。现有两个无功补偿方案：

【问题】
方案A：火电厂每台机组配置1台40Mvar的同步调相机，并配置2组60Mvar的并联电容器。风电场配置额定容量30\%的SVG（600Mvar）。
方案B：火电厂不配置调相机，仅配置4组60Mvar的并联电容器（总240Mvar），风电场配置额定容量20\%的SVG（400Mvar），并利用火电机组进相运行（功率因数0.95）吸收无功。

请依据GB 38755-2019和DL/T 5429-2009，对上述两个方案进行综合评估，并给出推荐方案及控制优先级链条。
  \item[A:] 现状诊断：该区域火电与风电并存，风电出力波动大，对无功电压控制要求高。根据GB 38755-2019第3.4条，无功补偿应分层分区平衡，并具备动态响应能力。DL/T 5429-2009第5.1.3条要求明确发电机功率因数及进相运行要求。当前问题在于如何选择无功补偿方案，确保电压稳定并兼顾经济性。

多维冲突分析：
1. 动态无功响应与稳态无功平衡的冲突：方案A配置调相机，可提供快速动态无功支撑（GB 38755-2019第5.6节要求电压稳定计算需考虑动态无功），但调相机投资高（约4000万元/台）。方案B仅靠SVG和电容器，SVG响应快但容量小，电容器响应慢，可能无法抑制风电波动引起的电压闪变。
2. 进相运行能力与机组安全的冲突：方案B要求火电机组进相运行（功率因数0.95），根据DL/T 5429-2009第5.1.3条，需校核发电机端部发热和定子端部铁芯温度，长期进相可能影响机组寿命。
3. 投资与运行成本的冲突：方案A总投资高（调相机+电容器+SVG约1.2亿元），但运行维护成本低。方案B投资低（约0.8亿元），但火电机组进相运行增加煤耗，且SVG容量不足时需弃风。

多方案比选：
- 方案A优势：动态无功支撑强，可满足GB 38755-2019第5.6节暂态电压稳定要求；调相机可提供短路容量，提升系统强度。劣势：投资高，占地大。
- 方案B优势：投资低，利用现有机组进相，无需新增设备。劣势：动态无功不足，风电波动时电压可能越限；火电机组进相运行风险高。

折中综合方案：推荐采用方案A，但优化配置：
1. 将调相机改为2台40Mvar（共80Mvar），配置在火电厂高压侧，满足GB 38755-2019第3.4条动态无功要求。
2. 风电场SVG容量提升至25\%（500Mvar），并配置快速电压控制功能。
3. 火电机组仍保留0.95（进相）能力，但仅作为紧急备用，不长期运行。

控制优先级链条：
第一优先级：调相机快速励磁响应，抑制电压骤降（GB 38755-2019第5.6节）。
第二优先级：风电场SVG动态无功输出，维持并网点电压（GB 38755-2019第3.4条）。
第三优先级：火电机组进相运行（紧急情况），吸收多余无功（DL/T 5429-2009第5.1.3条）。
第四优先级：并联电容器自动投切，调整稳态电压。
  \item[KW:] GB 38755-2019、DL/T 5429-2009、暂态电压稳定、分层分区、无功补偿
  \item[标准:] GB 38755-2019, DL/T 5429-2009
  \item[条款溯源:] GB 38755-2019 第3.2.1条, GB 38755-2019 第2.12条
  \item[主题:] 电压稳定
  \item[评分方式:] manual\_review
  \item[Rubric条款:] GB 38755-2019、DL/T 5429-2009、GB 38755-2019 第3.4条、DL/T 5429-2009 第5.1.3条、GB 38755-2019 第5.6节
  \item[Rubric判据:] 现状诊断：该区域火电与风电并存，风电出力波动大，对无功电压控制要求高、4条，无功补偿应分层分区平衡，并具备动态响应能力、3条要求明确发电机功率因数及进相运行要求、动态无功响应与稳态无功平衡的冲突：方案A配置调相机，可提供快速动态无功支撑（GB 38755-2019第5、6节要求电压稳定计算需考虑动态无功），但调相机投资高（约4000万元/台）
  \item[LJ判定:] 通过 (弱=0.60, 强=1.00, 差距=0.40)
\end{description}
\vspace{4pt}
\hrule
\vspace{6pt}

\end{document}
